\documentclass[10pt,a4paper]{article}
\usepackage[T1]{fontenc}
\usepackage[utf8]{inputenc}
\usepackage{lmodern}
\usepackage[left=22mm,right=22mm,top=18mm,bottom=18mm]{geometry}
\usepackage{parskip}
\usepackage{microtype}
\usepackage[hidelinks]{hyperref}

\pagestyle{empty}

\begin{document}

\begin{flushright}
Department of Intelligent Systems,\\
Faculty of Information Engineering,\\
Tokyo City University\\
1-28-1 Tamazutsumi, Setagaya-ku,\\
Tokyo 158-8557, Japan\\[4pt]
August 30, 2026
\end{flushright}

\noindent
Editors\\
\textit{Machine Learning and Knowledge Extraction} (MAKE)\\
MDPI

\vspace{1.5em}

\noindent
Dear Editors,

We are pleased to submit our manuscript entitled
\begin{quote}
\textbf{``Intrinsic-Dimension-Guided Bottleneck Selection for Autoencoders
and Variational Autoencoders with Active-Unit Diagnostics''}
\end{quote}
by Chika Obata, Mizuki Dai, and Kenya Jin'no, for consideration as a
Research Article in the \emph{Learning} section of \textit{Machine Learning
and Knowledge Extraction} (MAKE).

\textbf{Significance and context.}
The bottleneck dimension of autoencoders (AEs) and variational autoencoders
(VAEs) is one of their most fundamental design parameters, yet in practice
it is almost always set by convention or trial and error. Existing work has
measured the intrinsic dimension of data and of deep representations (e.g.,
Facco et al.\ 2017; Pope et al.\ 2021; Ansuini et al.\ 2019) or analyzed
posterior collapse in VAEs (Lucas et al.\ 2019), but stops short of
connecting these threads into a design procedure. Our manuscript closes
this gap with a conditional, diagnostic estimate--set--verify protocol: an
intrinsic-dimension estimate of the input data serves as a search-window
reference for the bottleneck dimension, and post-training verification
based on active units (AU) and reconstruction error diagnoses whether the
chosen dimension behaved as intended. The protocol is motivated
theoretically by a topological lower bound and a linear-Gaussian
rate--distortion surrogate, validated with multiple seeds on synthetic
manifolds with known ground truth and on MNIST and Fashion-MNIST, and
reduces the training runs of a full grid search by roughly 60--80\% on the
tested MNIST grid while containing the best observed linear-probe accuracy.
Importantly, the paper also quantifies where the approach reaches its
limits (natural images and convolutional VAEs), rather than claiming a
universal solution.

\textbf{Fit to the scope of MAKE.}
The paper extracts actionable knowledge from learned
representations---reading effective latent dimensionality from trained
models---and turns it into a transparent, reproducible design procedure
with pre-specified decision rules and a failure-mode taxonomy. This
combination of machine learning methodology and knowledge extraction,
together with the emphasis on reproducibility (all code, configurations,
and seed-specific logs are provided as an anonymized supplementary archive
for review, to be released publicly with a DOI upon acceptance), aligns
directly with the aims of MAKE.

A preliminary analysis of part of this study was presented at NOLTA~2025
(Obata \& Jin'no, ``Analysis of the Relationship Between Manifold Structure
and Latent Space Capacity of Input Data Using VAE''); it is cited in the
manuscript, which substantially extends it with the full protocol,
multi-seed validation, and applicability analysis. The manuscript has not
previously been submitted to MAKE or to any other MDPI journal. The use of
Claude Code (Anthropic; Claude Fable 5 model) for assistance in generating
experimental program code, English grammar checking, and proofreading is
disclosed with reasons in the
Acknowledgments section, and the authors take full responsibility for the
content. The authors declare no conflicts of interest.

We confirm that neither the manuscript nor any parts of its content are
currently under consideration for publication with or published in another
journal.

All authors have approved the manuscript and agree with its submission to
MAKE.

Thank you for considering our manuscript. We look forward to the
reviewers' feedback.

\vspace{0.8em}

\noindent
Sincerely,

\vspace{0.5em}

\noindent
Kenya Jin'no (corresponding author)\\
Department of Intelligent Systems, Faculty of Information Engineering,
Tokyo City University\\
E-mail: \href{mailto:kjinno@tcu.ac.jp}{kjinno@tcu.ac.jp};
Tel.: +81-3-5707-0104\\
On behalf of all authors: Chika Obata, Mizuki Dai, and Kenya Jin'no

\end{document}
